\documentclass[11pt]{article}
\usepackage[margin=1in]{geometry}
\usepackage{graphicx}
\usepackage{booktabs}
\usepackage{longtable}
\usepackage{array}
\usepackage{hyperref}
\usepackage{float}
\usepackage{setspace}
\setstretch{1.08}

\title{CC-SPR as a Curvature-Corrected Sparse Extension of Harmonic Persistent Homology:\\Public CLL Substitute Study}
\author{CC-SPR Reproducible Pipeline}
\date{March 17, 2026}

\begin{document}
\maketitle

\begin{abstract}
We implemented a reproducible CC-SPR pipeline and evaluated it on public CLL substitutes after access limitations blocked the originally requested CLL datasets. The substitute datasets were GEO GSE161711 (bulk RNA-seq) and GSE165087 (single-cell RNA-seq), selected to retain CLL and longitudinal-treatment relevance. The pipeline includes real data ingestion, preprocessing, Euclidean and Ricci geometry options, sparse cycle representatives, TIP stability, full ablation grids, and figure/table generation. This document summarizes methods, results, limitations, and next steps under a local-scale execution profile.
\end{abstract}

\section{Context and Substitution Rationale}
The original project framing was preserved: CC-SPR extends harmonic PH-style omics probing by combining Ricci-corrected geometry with sparse cycle representatives and TIP feature stability.

Because the original CLL venetoclax matrix and controlled-access CLL/RS scRNA cohort were unavailable in this environment, we used public substitutes:
\begin{itemize}
\item Bulk substitute: \textbf{GSE161711}
\item scRNA substitute: \textbf{GSE165087}
\end{itemize}

These substitutes are explicitly treated as replacements chosen for CLL modality alignment (bulk + single-cell), not as exact cohort matches.

\section{Methods}
For each split, the pipeline performs:
\begin{enumerate}
\item Variable-feature selection.
\item Distance construction (Euclidean or Ricci-corrected).
\item Persistent homology computation and dominant H1 cycle extraction.
\item Cycle representative optimization (L2 or elastic).
\item Projection to feature scores and TIP bootstrap aggregation.
\item Weighted-F1 evaluation with logistic regression.
\end{enumerate}

\noindent Compared model families:
\begin{itemize}
\item \texttt{standard}
\item \texttt{harmonic} (Euclidean + L2)
\item \texttt{eu\_sparse} (Euclidean + elastic)
\item \texttt{ccspr} (Ricci + elastic)
\end{itemize}

\noindent Full ablation grid retained:
\begin{itemize}
\item distance mode \{euclidean, ricci\}
\item Ricci iterations \{0, 2, 5, 10\}
\item k \{5, 10, 15\}
\item lambda \{1e-8, 1e-7, 1e-6, 1e-5, 1e-4\}
\item top-K \{10, 20, 50\}
\item projection normalize \{none, std\}
\end{itemize}

\section{Data Construction}
\textbf{GSE161711 bulk} was converted to:
\begin{itemize}
\item \texttt{data/cll\_venetoclax/processed\_matrix.tsv}
\item \texttt{data/cll\_venetoclax/sample\_metadata.tsv}
\end{itemize}

\textbf{GSE165087 scRNA} was converted to:
\begin{itemize}
\item \texttt{data/cll\_rs\_scrna/raw/raw\_combined.h5ad}
\item \texttt{data/cll\_rs\_scrna/data.h5ad}
\item \texttt{data/cll\_rs\_scrna/sample\_metadata.tsv}
\end{itemize}

Raw and processed outputs were kept separate.

\section{Results}
\subsection{Main Benchmark}
\begin{table}[H]
\centering
\caption{Weighted F1 means and 95\% CI (local-scale runs).}
\begin{tabular}{llll}
\toprule
Dataset & Model & Mean F1 & 95\% CI \\
\midrule
GSE161711 bulk & standard & 0.8961 & 0.0404 \\
GSE161711 bulk & harmonic & 0.4804 & 0.1107 \\
GSE161711 bulk & eu\_sparse & 0.5885 & 0.0931 \\
GSE161711 bulk & ccspr & 0.4818 & 0.0028 \\
GSE165087 scRNA & standard & 0.7819 & 0.0994 \\
GSE165087 scRNA & harmonic & 0.5214 & 0.1554 \\
GSE165087 scRNA & eu\_sparse & 0.4704 & 0.0466 \\
GSE165087 scRNA & ccspr & 0.4734 & 0.1438 \\
\bottomrule
\end{tabular}
\end{table}

\subsection{Ablation Summary}
Observed best mean settings in these local runs:
\begin{itemize}
\item GSE161711 bulk: strong values at \texttt{k=15}, \texttt{top\_k=50}, \texttt{normalize=std}, \texttt{lambda=1e-4}.
\item GSE165087 scRNA: strong values at \texttt{iters=10}, \texttt{k=5}, \texttt{top\_k=50}, \texttt{lambda=1e-7}.
\end{itemize}

\subsection{Figures}
\begin{figure}[H]
\centering
\includegraphics[width=0.88\textwidth]{../results/figures/main_benchmark_public_substitutes.png}
\caption{Main benchmark summary for public CLL substitutes.}
\end{figure}

\begin{figure}[H]
\centering
\includegraphics[width=0.95\textwidth]{../results/figures/ablation_panels_public_substitutes.png}
\caption{Ablation panels across both substitute datasets.}
\end{figure}

\begin{figure}[H]
\centering
\includegraphics[width=0.48\textwidth]{../results/figures/cll_venetoclax_tip_eu_vs_ricci.png}
\includegraphics[width=0.48\textwidth]{../results/figures/cll_rs_scrna_tip_eu_vs_ricci.png}
\caption{TIP comparison (Euclidean vs Ricci) for bulk and scRNA substitutes.}
\end{figure}

\section{Limitations}
\begin{enumerate}
\item The substitutes are biologically aligned but are not the original requested cohorts.
\item Current scRNA results are local-scale for feasibility; full-scale reruns are recommended.
\item Current split/boot settings prioritize reproducible local completion over maximum statistical power.
\end{enumerate}

\section{Future Work}
\begin{enumerate}
\item Full-scale reruns on larger cell budgets and higher bootstrap counts.
\item Patient/timepoint-level endpoint expansion.
\item Explicit R/maTilDA parity benchmarking in final tables.
\item Multi-seed repeated evaluation and nested cross-validation.
\end{enumerate}

\section{Reproducibility}
Core commands:
\begin{verbatim}
source ~/.venvs/lesegenv/bin/activate
export PYTHONPATH=src
export MPLCONFIGDIR=/tmp/matplotlib
export NUMBA_DISABLE_COVERAGE=1
export NUMBA_CACHE_DIR=/tmp/numba_cache

python3 scripts/download_cll_venetoclax.py --root data/cll_venetoclax
python3 scripts/download_cll_rs_scrna.py --root data/cll_rs_scrna --max-cells-raw 40000 --max-cells-processed 12000 --seed 42
python3 experiments/run_cll_venetoclax.py --config configs/cll_venetoclax.yaml
python3 experiments/run_cll_rs_scrna.py --config configs/cll_rs_scrna.yaml
python3 experiments/aggregate_public_substitutes.py
\end{verbatim}

\end{document}
